An den Rändern steht eine $1$; dazwischen entsteht jede Zahl als Summe
der beiden Zahlen schräg darüber. Diese Lesefassung erklärt die
Binomialaussagen aus Abschnitt~5.3 des Leitfadens zu Band~10 mit
vollständigen Beweisen. Wege durch das Dreieck verbinden die einzelnen
Themen; die Zuordnung zu den Leitfadensätzen steht in Abschnitt~\ref{bin:map}.

\textbf{Voraussetzungen und Schreibweise.}
Wir verwenden $\N=\{0,1,2,\ldots\}$, Induktion, Rekursion und die
Rechengesetze der natürlichen Zahlen. Es gilt $a^0=1$, auch für $a=0$;
eine leere Summe ist $0$. Differenzen wie $n-k$ bilden wir nur für $k\le n$.
Die arithmetischen Beweise bleiben in $\N$; Wege und Auswahlen ergänzen
sie durch elementares Zählen.

\textbf{Mit dem Heft arbeiten.}
Nutze die Karoseite für eigene Wege, Rechenschritte und Fragen.
Acht Aufgaben mit Arbeitsraum und Lösungen stehen im separaten
\href{_B10-binomial-exercises.pdf\#binomial.exercises}{Übungsheft}.

\medskip\tableofcontents
\clearpage

\section{Das Dreieck entsteht Zeile für Zeile}
\label{bin:construction}

Die ersten sieben Zeilen sehen so aus:
\begin{center}
\PascalDreieck{1}
\end{center}
Die Zeilenzählung beginnt bei $n=0$. Innerhalb der Zeile zählen wir
von links ebenfalls ab $0$. Der Eintrag an Stelle $k$ in Zeile $n$
heißt \emph{Binomialkoeffizient} und wird $\binom nk$ geschrieben,
gesprochen: \glqq $n$ über $k$\grqq. Beispielsweise ist
\[
 \binom42=6=3+3=\binom31+\binom32.
\]
Die untere Zahl ist ein Index, kein Nenner. Das Symbol bezeichnet
einen einzelnen Eintrag des Dreiecks.

\subsection*{Eine Regel, die auch die Ränder erfasst}

Für Rechnungen ist es praktisch, jede Zeile rechts durch Nullen
fortzusetzen. So lautet etwa die ganze Zeile $4$:
$1,4,6,4,1,0,0,\ldots$. Die Nullen werden im Dreieck nicht mitgezeichnet.

\begin{definition}[Pascal-Rekursion]
Für $n,k\in\N$ setzen wir
\begin{equation}\tag{P}\label{bin:pascal}
 \begin{aligned}
 \binom00&=1,&\binom0{k+1}&=0,\\
 \binom{n+1}0&=1,&
 \binom{n+1}{k+1}&=\binom nk+\binom n{k+1}.
 \end{aligned}
\end{equation}
\end{definition}
Die Regel für $k+1$ addiert genau die beiden Einträge schräg darüber.
Zum Beispiel entsteht die nächste Zeile aus $1,3,3,1,0,\ldots$ durch
\[
 1,\quad 1+3,\quad 3+3,\quad 3+1,\quad 1+0,\quad 0+0,\ldots.
\]
So erhält man $1,4,6,4,1,0,\ldots$.

\textbf{Warum ist damit jeder Eintrag eindeutig festgelegt?}
Die Anfangszeile ist vollständig bekannt. Ist eine Zeile bekannt,
so bestimmen die beiden Regeln jede Stelle der nächsten Zeile durch
eine einzige natürliche Zahl. Induktion zeigt daher, dass zwei
Zahlensysteme mit diesen Regeln in jeder Zeile übereinstimmen.
Auch die Existenz folgt aus dem Rekursionssatz von Band~10:
Man betrachtet eine Zeile als Abbildung $f\colon\N\to\N$ und bildet
daraus die neue Zeile $Tf$ durch
\[
 (Tf)(0)=1,\qquad (Tf)(k+1)=f(k)+f(k+1).
\]
Der Operator $T$ führt Abbildungen $\N\to\N$ wieder in solche
Abbildungen über. Seine wiederholte Anwendung auf die Anfangszeile
liefert die gesuchte Folge von Zeilen.

\subsection{Was die Zahlen zählen: Wege durch das Dreieck}
\label{bin:paths}

Wir starten an der Spitze $(0,0)$. Ein Schritt führt schräg nach links
unten ($L$) oder rechts unten ($R$). Die Stelle zählt innerhalb jeder
Zeile ab $0$; deshalb gilt
\[
 L:(n,k)\longmapsto(n+1,k),\qquad
 R:(n,k)\longmapsto(n+1,k+1).
\]
Nach $n$ Schritten mit genau $k$ Rechtsschritten endet der Weg bei
$(n,k)$. Der obere Index zählt also \emph{alle Schritte}, der untere
\emph{die Rechtsschritte}. Die Zeichnung zeigt den Weg $RLLR$ nach $(4,2)$.
\begin{center}\PascalWegeStart\end{center}

\begin{satz}[Wegeinterpretation]
Die Anzahl der Wege von der Spitze nach $(n,k)$ ist $\binom nk$.
Für $k>n$ gibt es keinen solchen Weg; die Anzahl ist $0$.
\end{satz}
\begin{proof}
Nennen wir die Anzahl zunächst $W(n,k)$. An der Spitze gibt es genau
einen Weg ohne Schritte: $W(0,0)=1$. Ohne Schritte erreicht man keine
positive Stelle, also $W(0,k+1)=0$. Zum linken Rand führt nur der Weg
aus lauter $L$, daher $W(n+1,0)=1$.

Ein Weg nach $(n+1,k+1)$ endet entweder mit $R$ von $(n,k)$ oder
mit $L$ von $(n,k+1)$. Entfernen des letzten Schritts und erneutes
Anhängen sind zueinander umkehrbar. Die beiden Klassen sind disjunkt
und erfassen alle Wege. Deshalb
\[
 W(n+1,k+1)=W(n,k)+W(n,k+1).
\]
Anfangswerte und Rekursion stimmen mit \eqref{bin:pascal} überein.
Die dort bewiesene Eindeutigkeit liefert $W(n,k)=\binom nk$.
Auch außerhalb des Dreiecks passt dies: Mehr als $n$ Rechtsschritte
in insgesamt $n$ Schritten sind unmöglich.
\end{proof}

Zum Beispiel führen genau die sechs Schrittfolgen
\[
 LLRR,\ LRLR,\ LRRL,\ RLLR,\ RLRL,\ RRLL
\]
nach $(4,2)$. Somit zählt die $6$ an dieser Stelle sechs verschiedene
Wege; die Zeichnung hebt einen davon hervor.

\subsection{Die Indizes merken: einer weniger, einer gleich}
\label{bin:indices}

Wir benennen den \emph{gesuchten Eintrag} zuerst: Er steht in der neuen
Zeile $n+1$ an der Stelle $k$. Für $k\ge1$ lautet die Pascal-Regel dann
\begin{equation}\tag{P'}\label{bin:pascal-k}
 \boxed{\displaystyle
 \binom{n+1}k=\binom n{k-1}+\binom nk.}
\end{equation}
Beide Eltern stehen eine Zeile höher, also in Zeile $n$.
Links oben steht der Index $k-1$, rechts oben der Index $k$:
\begin{center}\PascalEltern\end{center}

Der Index $k$ bezeichnet keine feste senkrechte Bildspalte.
Die längere neue Zeile wird beim Zentrieren um einen halben Abstand
nach links versetzt. Deshalb liegt ihr Eintrag $k$ zwischen den alten
Stellen $k-1$ und $k$.

\textbf{Mit Wegen merken.} Ist der letzte Schritt ein $R$, enthält er
schon einen der insgesamt $k$ Rechtsschritte: Vorher waren es $k-1$.
Ist der letzte Schritt ein $L$, bleibt ihre Anzahl $k$.
So entstehen genau die beiden Summanden in \eqref{bin:pascal-k}.

\medskip
\noindent\fcolorbox{pascalblue!55}{pascalpale}{%
\begin{minipage}{\dimexpr\linewidth-2\fboxsep-2\fboxrule\relax}
\textbf{Merksatz:}\quad Oben eine Zeile zurück;\\
unten einmal weniger, einmal gleich.
\end{minipage}}

Am linken Rand $k=0$ gilt direkt $\binom{n+1}0=1$; ein Index $-1$
wird nicht benötigt. Die Regeln (P) und (P') sind dieselbe Aussage:
In (P) heißt die Zielstelle $k+1$, in (P') heißt sie $k$.

\section{Randwerte, Nullen und Spiegelung}
\label{bin:basic}

Im Wegenetz sind die Randwerte sofort verständlich: Am linken Rand
führt jeder Schritt nach links, am rechten jeder nach rechts. Beide
Wege sind eindeutig. Jede Stelle $0\le k\le n$ ist erreichbar, etwa
mit zuerst $k$ Schritten $R$ und dann $n-k$ Schritten $L$.
Für $\binom n1$ kann der einzige Rechtsschritt an jeder der $n$
Positionen stehen. Der folgende Induktionsbeweis bestätigt alle Aussagen.

\begin{satz}[Randwerte und positive Einträge]
Für alle $n,k\in\N$ gilt
\begin{equation}\tag{G}\label{bin:boundary}
 \binom n0=\binom nn=1,\qquad
 \binom nk=0\ \Longleftrightarrow\ k>n,\qquad
 \binom n1=n.
\end{equation}
Insbesondere sind alle Einträge mit $0\le k\le n$ positiv.
\end{satz}

\begin{proof}
Wir beweisen zuerst Randwerte, Positivität und Nullen gemeinsam durch
Induktion über $n$. In Zeile $0$ steht an Stelle $0$ die Zahl $1$;
an allen späteren Stellen steht $0$. Damit gelten diese Aussagen am Anfang.

Sie seien für Zeile $n$ richtig. Der linke Rand der nächsten Zeile ist
nach Definition $1$. Am rechten Rand liefert die Rekursion
\[
 \binom{n+1}{n+1}=\binom nn+\binom n{n+1}=1+0=1.
\]
Für $1\le k\le n+1$ gilt
\[
 \binom{n+1}k=\binom n{k-1}+\binom nk.
\]
Der erste Summand ist positiv, weil $0\le k-1\le n$; der zweite ist
eine natürliche Zahl. Daher ist die Summe positiv.
Für $k>n+1$ liegen dagegen sowohl $k-1$ als auch $k$ rechts außerhalb
der alten Zeile. Beide Summanden sind $0$. Das beweist zugleich,
dass genau die Stellen außerhalb des Dreiecks verschwinden.

Für die letzte Formel beginnen wir mit $\binom01=0$. Aus
$\binom n1=n$ folgt
\[
 \binom{n+1}1=\binom n0+\binom n1=1+n=n+1.
\]
Eine zweite Induktion liefert also $\binom n1=n$ für alle $n$.
\end{proof}

\begin{satz}[Symmetrie]
Für $0\le k\le n$ gilt
\begin{equation}\tag{S}\label{bin:symmetry}
 \binom nk=\binom n{n-k}.
\end{equation}
\end{satz}

\Beweisidee Vertausche in jedem Weg $L$ und $R$. Ein Weg mit $k$
Rechtsschritten wird zu einem mit $n-k$ Rechtsschritten; erneutes
Vertauschen gibt den ursprünglichen Weg zurück. Diese eindeutige
Paarung erklärt die Symmetrie. Arithmetisch folgt sie so:

\begin{proof}
Die Zeile $0$ besteht im Dreieck nur aus einer $1$ und ist symmetrisch.
Sei die Symmetrie für Zeile $n$ bewiesen. In Zeile $n+1$ stimmen
die beiden Ränder überein, weil dort jeweils $1$ steht.
Für einen inneren Index $1\le k\le n$ gilt
\[
 \begin{aligned}
 \binom{n+1}k
 &=\binom n{k-1}+\binom nk\\
 &=\binom n{n+1-k}+\binom n{n-k}\\
 &=\binom{n+1}{n+1-k}.
 \end{aligned}
\]
Die mittlere Gleichheit verwendet zweimal die Induktionsannahme,
die letzte Gleichheit die Pascal-Rekursion und die Vertauschbarkeit
der Addition. Damit ist auch Zeile $n+1$ symmetrisch.
\end{proof}

\begin{beispiel}[Was sich sofort ablesen lässt]
In Zeile $6$ stehen $1,6,15,20,15,6,1$. Insbesondere sind
$\binom62=\binom64=15$ und $\binom65=\binom61=6$.
Dagegen ist $\binom67=0$. Eine Stelle rechts vom Dreieck ist
wohldefiniert; sie hat den Wert $0$.
\end{beispiel}

\subsection{Ein Weg ist zugleich eine Auswahl}
\label{bin:selection}

Ein Weg ist eindeutig bestimmt, sobald die Positionen seiner
Rechtsschritte feststehen. Beispielsweise gehört zur Schrittfolge
$RLLR$ die Auswahl $\{1,4\}$ aus den vier Schrittpositionen.
Umgekehrt macht jede Auswahl von $k$ Positionen aus $n$ einen Weg:
An den gewählten Stellen gehen wir nach rechts, sonst nach links.
Beide Zuordnungen heben einander auf. Daher zählt $\binom nk$ auch
die Auswahlen von $k$ aus $n$ verschiedenen Gegenständen, ohne
Reihenfolge und ohne Wiederholung.

Für $n=4$, $k=2$ entsprechen die sechs Wege den sechs Auswahlen
\[
 \{1,2\},\ \{1,3\},\ \{1,4\},\ \{2,3\},\ \{2,4\},\ \{3,4\}.
\]
Die Pascal-Zerlegung nach dem letzten Schritt wird dabei zur Zerlegung
nach der letzten Position: Sie ist ausgewählt ($R$) oder nicht ($L$).
Damit beschreiben Wege und Auswahlen dieselben Zahlen.

\section{Ein einzelner Eintrag ohne die vorherigen Zeilen}
\label{bin:factorials}

Die Rekursion ist bequem, wenn man eine ganze Zeile benötigt.
Für einen einzelnen Eintrag weit unten im Dreieck wäre es umständlich,
alle vorherigen Zeilen aufzubauen. Hier helfen die Fakultäten.

Wir erinnern an
\[
 0!=1,\qquad (n+1)!=(n+1)n!.
\]
Insbesondere sind $1!=1$, $2!=2$, $3!=6$, $4!=24$.
Jede Fakultät ist positiv: Das gilt für $0!=1$; ist $n!>0$,
so ist auch $(n+1)n!>0$. Damit ist die Behauptung durch Induktion bewiesen.

\begin{satz}[Fakultätenformel]
Für $0\le k\le n$ gilt
\begin{equation}\tag{F}\label{bin:factorial}
 \boxed{\displaystyle k!(n-k)!\binom nk=n!.}
\end{equation}
Der Binomialkoeffizient ist die einzige natürliche Zahl, die diese
Gleichung erfüllt.
\end{satz}

\Beweisidee Ein Weg nach $(n,k)$ ist eine Anordnung von $k$ Zeichen $R$
und $n-k$ Zeichen $L$. Beschriften wir alle Zeichen einzeln, gibt es
$n!$ Anordnungen: Für die Plätze bestehen nacheinander $n,n-1,\ldots,1$
Wahlmöglichkeiten. Jeder unbeschriftete Weg entsteht dabei genau
$k!(n-k)!$-mal, durch Vertauschen der $R$-Beschriftungen und der
$L$-Beschriftungen. So erklärt das Wegemodell bereits die Gleichung (F).
Für null Zeichen gibt es eine leere Anordnung. Es folgt der arithmetische
Beweis ohne vorausgesetzte Division.

\begin{proof}
Wir führen Induktion über $n$, jeweils für alle $k\le n$.
Für $n=0$ ist nur $k=0$ möglich, und es gilt $0!0!\binom00=1=0!$.

Die Formel gelte nun in Zeile $n$. In der nächsten Zeile sind die
Randfälle $k=0$ und $k=n+1$ sofort richtig, weil die Binomialkoeffizienten
an beiden Rändern $1$ sind. Sei also $1\le k\le n$.
Mit der Rekursion erhalten wir
\[
 \begin{aligned}
 &k!(n+1-k)!\binom{n+1}k\\
 &\quad=k!(n+1-k)!\binom n{k-1}
       +k!(n+1-k)!\binom nk.
 \end{aligned}
\]
Für den ersten Summanden benutzen wir $k!=k(k-1)!$ und die
Induktionsannahme am Index $k-1$:
\[
 \begin{aligned}
 &k!(n+1-k)!\binom n{k-1}\\
 &\quad=k\bigl((k-1)!(n-(k-1))!\binom n{k-1}\bigr)\\
 &\quad=kn!.
 \end{aligned}
\]
Für den zweiten benutzen wir $(n+1-k)!=(n+1-k)(n-k)!$ und die
Induktionsannahme am Index $k$:
\[
 \begin{aligned}
 &k!(n+1-k)!\binom nk\\
 &\quad=(n+1-k)\bigl(k!(n-k)!\binom nk\bigr)\\
 &\quad=(n+1-k)n!.
 \end{aligned}
\]
Zusammen ergibt dies
\[
 kn!+(n+1-k)n!=(n+1)n!=(n+1)!.
\]
Damit ist die Formel bewiesen. Schließlich ist $k!(n-k)!$ positiv.
Erfüllen zwei natürliche Zahlen die Gleichung \eqref{bin:factorial},
kann man diesen gemeinsamen positiven Faktor kürzen; die Zahlen sind gleich.
\end{proof}

Die vertraute Schreibweise
\[
 \binom nk=\frac{n!}{k!(n-k)!}\qquad(0\le k\le n)
\]
bezeichnet somit eine \emph{exakte Division}: Der Zähler ist ein
natürliches Vielfaches des Nenners. Die Existenz des Ergebnisses in
$\N$ wurde durch die Rekursion gesichert, bevor wir einen Bruch schreiben.
Für $k>n$ verwenden wir weiterhin den Wert $0$ aus \eqref{bin:boundary};
eine Fakultät von $n-k$ wird dann nicht gebildet.

\begin{beispiel}[Zwei Plätze unter zehn]
Für $\binom{10}2$ liefert die Fakultätenformel
\[
 2!\,8!\binom{10}2=10!=10\cdot9\cdot8!.
\]
Nach Kürzen von $8!$ bleibt $2\binom{10}2=90$, also
$\binom{10}2=45$. Nach der Symmetrie ist auch $\binom{10}8=45$.
Man muss weder die zehnte Zeile noch große Fakultäten ausrechnen.
\end{beispiel}

\section{Warum das Dreieck beim Potenzieren erscheint}
\label{bin:theorem}

Beim Ausmultiplizieren von $(a+b)^2$ entstehen die drei Terme
$b^2+2ab+a^2$, bei $(a+b)^3$ die vier Terme
$b^3+3ab^2+3a^2b+a^3$. Die Koeffizienten sind die Zeilen $2$ und $3$.
Wir ordnen die Terme hier nach aufsteigenden Potenzen von $a$.

\begin{satz}[Binomischer Lehrsatz]
Für $a,b,n\in\N$ gilt
\begin{equation}\tag{B}\label{bin:binomial-theorem}
 \boxed{\displaystyle
 (a+b)^n=\sum_{k=0}^n\binom nk a^k b^{n-k}.}
\end{equation}
\end{satz}

\Beweisidee Beim Ausmultiplizieren der $n$ Faktoren $(a+b)$ wählen wir
in jedem Faktor $a$ oder $b$. Lies $a$ als Rechtsschritt und $b$ als
Linksschritt. Jede Schrittfolge liefert genau ein Produkt. Die
$\binom nk$ Wege nach $(n,k)$ liefern alle $a^kb^{n-k}$; nach Endpunkten
zusammengefasst ergibt sich (B). Auch bei $a=0$ oder $b=0$ bleibt dies
richtig, mit der vereinbarten Potenzkonvention. Die folgende Induktion
führt das Zusammenfassen der Terme vollständig aus.

\begin{proof}
Wir halten $a,b$ fest und führen Induktion über $n$.
Für $n=0$ ist die linke Seite $1$. Rechts steht nur
$\binom00a^0b^0=1$. Dies gilt auch, wenn $a$ oder $b$ null ist.

Sei die Formel für $n$ bewiesen. Dann ergeben das Distributivgesetz
und die Potenzregeln
\[
 \begin{aligned}
 (a+b)^{n+1}
 &=(a+b)\sum_{k=0}^n\binom nk a^k b^{n-k}\\
 &=\sum_{k=0}^n\binom nk a^{k+1}b^{n-k}\\
 &\quad+\sum_{k=0}^n\binom nk a^k b^{n+1-k}.
 \end{aligned}
\]
In der ersten Summe setzen wir $j=k+1$. Aus den Grenzen $0\le k\le n$
werden $1\le j\le n+1$, und die erste Summe wird zu
\[
 \sum_{j=1}^{n+1}\binom n{j-1}a^j b^{n+1-j}.
\]
Nun benennen wir $j$ wieder in $k$ um. Für $1\le k\le n$ enthalten
beide Summen denselben Faktor $a^k b^{n+1-k}$.
Der Term $k=0$ kommt nur aus der zweiten, der Term $k=n+1$ nur
aus der ersten Summe. Mit den Randwerten $\binom n0=\binom nn=1$
erhalten wir daher
\[
 \begin{aligned}
 (a+b)^{n+1}
 &=b^{n+1}+a^{n+1}\\
 &\quad+\sum_{k=1}^n\left(\binom n{k-1}+\binom nk\right)
        a^k b^{n+1-k}\\
 &=b^{n+1}+a^{n+1}\\
 &\quad+\sum_{k=1}^n\binom{n+1}k a^k b^{n+1-k}\\
 &=\sum_{k=0}^{n+1}\binom{n+1}k a^k b^{n+1-k}.
 \end{aligned}
\]
Beim Zusammenfassen der Koeffizienten wird die Pascal-Rekursion
verwendet. In der letzten Zeile werden
die beiden Randterme wieder in die Summe aufgenommen. Ist $n=0$,
ist die innere Summe leer; dieselbe Rechnung liefert dann $b+a$.
\end{proof}

\begin{beispiel}[Die vierte Potenz]
Die Zeile $1,4,6,4,1$ liefert
\[
 (a+b)^4=b^4+4ab^3+6a^2b^2+4a^3b+a^4.
\]
Für $a=2$ und $b=3$ wird daraus
\[
 5^4=81+216+216+96+16=625.
\]
Die sechs Wege nach $(4,2)$ entsprechen genau den sechs Produkten
mit zweimal $a$ und zweimal $b$. Deshalb steht vor $a^2b^2$ die Zahl $6$.
\end{beispiel}

\section{Zeilen und Diagonalen summieren}
\label{bin:sums}

Im Wegemodell erklären wir Summen durch Ordnen der Wege:
nach ihrem Endpunkt für Zeilensummen und nach dem letzten
Rechtsschritt für Diagonalensummen.

\begin{satz}[Zeilensumme]
Für jedes $n\in\N$ gilt
\begin{equation}\tag{Z}\label{bin:rowsum}
 \sum_{k=0}^n\binom nk=2^n.
\end{equation}
\end{satz}
\begin{proof}
Bei jedem Schritt gibt es unabhängig zwei Richtungen. Daher gibt es
$2^n$ Wege der Länge $n$, bei $n=0$ genau einen leeren Weg.
Jeder endet an genau einer Stelle $k$ der Zeile $n$; diese Klasse hat
$\binom nk$ Wege. Addition über die disjunkten Endpunktklassen beweist
die Formel. Arithmetisch erhält man dieselbe Gleichung aus
\eqref{bin:binomial-theorem} mit $a=b=1$.
\end{proof}

Beispielsweise zählt $1+4+6+4+1=16=2^4$ alle Wege mit vier Schritten.

\subsection*{Die Summe entlang einer Diagonale}

Halten wir den unteren Index $r$ fest und erhöhen nur den oberen,
lesen wir eine schräge Diagonale. Bei $r=2$ beginnt sie mit
$1,3,6,10,15,\ldots$. Eine endliche Anfangssumme steht eine Zeile
tiefer in der nächsten Diagonale:
\begin{center}
\PascalWegeDiagonal
\end{center}
Die markierten Einträge zeigen
\[
 \binom22+\binom32+\binom42+\binom52
 =1+3+6+10=20=\binom63.
\]

\begin{satz}[Diagonalensumme]
Für alle $n,r\in\N$ gilt
\begin{equation}\tag{D}\label{bin:diagonal}
 \boxed{\displaystyle
 \sum_{i=0}^n\binom{i+r}r=\binom{n+r+1}{r+1}.}
\end{equation}
\end{satz}

\begin{proof}[Beweis mit Wegen]
Ein Weg zum Ziel $(n+r+1,r+1)$ hat $r+1$ Rechtsschritte und $n$
Linksschritte. Sein letzter Rechtsschritt folgt nach genau $i$
Linksschritten und $r$ Rechtsschritten, für ein eindeutiges $i$ mit
$0\le i\le n$. Der Weg davor endet also bei $(i+r,r)$ und kann auf
$\binom{i+r}r$ Arten gewählt werden. Danach folgen der letzte Schritt
$R$ und zwingend $n-i$ Schritte $L$. Umgekehrt ergibt jeder solche
Anfangsweg genau einen vollständigen Weg dieser Klasse.

Die Klassen für $i=0,\ldots,n$ sind disjunkt und erfassen alle Wege
zum Ziel. Ihre Anzahlen addieren sich deshalb zu
$\binom{n+r+1}{r+1}$. Im Bild ist $r=2$, $n=3$: Die Zahlen
$1,3,6,10$ zählen die möglichen Anfangswege vor dem letzten $R$.
\end{proof}

\begin{proof}[Arithmetischer Beweis]
Wir halten $r$ fest und führen Induktion über $n$.
Für $n=0$ lautet die Gleichung
$\binom rr=\binom{r+1}{r+1}$; beide Seiten sind $1$.
Die Formel gelte für $n$. Für die um einen Summanden längere Summe gilt
\[
 \begin{aligned}
 \sum_{i=0}^{n+1}\binom{i+r}r
 &=\left(\sum_{i=0}^n\binom{i+r}r\right)+\binom{n+r+1}r\\
 &=\binom{n+r+1}{r+1}+\binom{n+r+1}r\\
 &=\binom{n+r+2}{r+1}.
 \end{aligned}
\]
Die zweite Gleichheit ist die Induktionsannahme, die dritte die
Pascal-Rekursion. Dies ist genau die Behauptung für $n+1$.
\end{proof}

\begin{beispiel}[Die Summe der ersten natürlichen Zahlen]
Für $r=1$ und $\binom{i+1}1=i+1$ liefert \eqref{bin:diagonal}
\[
 1+2+\cdots+(n+1)=\binom{n+2}2.
\]
Aus der Fakultätenformel folgt durch Kürzen von $n!$
\[
 2\binom{n+2}2=(n+1)(n+2).
\]
Damit erhalten wir die natürliche Gleichung
\[
 2\sum_{i=0}^n(i+1)=(n+1)(n+2).
\]
Bei $n=3$ ist dies $2(1+2+3+4)=4\cdot5$.
Für $r=0$ besagt dieselbe Diagonalensummenformel einfach:
$n+1$ Einsen ergeben $\binom{n+1}1=n+1$.
\end{beispiel}

\section{Zwei Zeilen zusammenfügen: Vandermonde}
\label{bin:vandermonde-section}

Die Pascal-Regel ordnet Wege nach ihrem letzten Schritt.
Nun wählen wir eine beliebige Zwischenzeile: Jeder Weg trifft sie an
genau einer Stelle. An jeder Station kombinieren wir Hinwege und
Weiterwege; anschließend addieren wir über alle Stationen.

\begin{satz}[Vandermondesche Identität]
Für alle $m,n,r\in\N$ gilt
\begin{equation}\tag{V}\label{bin:vandermonde}
 \boxed{\displaystyle
 \sum_{k=0}^r\binom mk\binom n{r-k}=\binom{m+n}r.}
\end{equation}
\end{satz}

\begin{center}\PascalWegeVandermonde\end{center}
Im Bild führen $2$ Hinwege zur Station B und $3$ Weiterwege von B
zum Ziel. Jeder Hinweg passt zu jedem Weiterweg: Über B führen
$2\cdot3=6$ Wege. Über A und C kommen $1\cdot1$ und $1\cdot3$ hinzu.
Insgesamt sind es $10=\binom53$ Wege.

\begin{proof}[Beweis mit Wegen]
Sei zunächst $r\le m+n$. Ein Weg nach $(m+n,r)$ hat insgesamt
$m+n$ Schritte, darunter $r$ Rechtsschritte. Nach den ersten $m$
Schritten steht er an genau einer Stelle $(m,k)$, wobei $0\le k\le r$.
Den Hinweg kann man auf $\binom mk$ Arten wählen.

Für den Weiterweg bleiben $n$ Schritte, davon $r-k$ nach rechts.
Beginnt man die Schrittzählung an der Zwischenstation neu, gibt es
also $\binom n{r-k}$ Weiterwege. Diese Zahl zählt \emph{ab der Station};
sie ist nicht der globale Eintrag an einem späteren Punkt des Dreiecks.
Jeder Hinweg lässt sich mit jedem Weiterweg verbinden. Umgekehrt
gewinnt man beide Teile eindeutig zurück, indem man nach dem $m$-ten
Schritt trennt. Somit gibt es genau
$\binom mk\binom n{r-k}$ Wege über diese Station.

Die Stationsklassen sind disjunkt und vollständig. Addition beweist
(V). Unmögliche Stationen oder Weiterwege tragen $0$ bei. Ist
$r>m+n$, ist das Ziel unerreichbar. Auch jeder Summand ist dann $0$:
Aus $k\le m$ und $r-k\le n$ würde $r\le m+n$ folgen. Damit gilt die
Gleichung auch außerhalb des Dreiecks.
\end{proof}

\textbf{Dasselbe als Kugelauswahl.}
Teile $m+n$ unterscheidbare Kugeln in Gruppen mit $m$ und $n$ Kugeln.
Jede Auswahl von $r$ Kugeln enthält eindeutig $k$ aus der ersten und
$r-k$ aus der zweiten Gruppe. Für festes $k$ werden die Möglichkeiten
multipliziert, über die möglichen $k$ addiert. Im Wegemodell sind die
ausgewählten Positionen gerade die Rechtsschritte in den beiden Wegteilen.

\begin{proof}[Arithmetischer Beweis]
Wir halten $m$ fest und setzen
\[
 V_n(r):=\sum_{k=0}^r\binom mk\binom n{r-k}.
\]
Wir führen Induktion \emph{nur über $n$}, jeweils für alle $r\in\N$.
Bei $n=0$ ist $\binom0{r-k}$ nur für $k=r$ ungleich $0$ und dort $1$.
Daher $V_0(r)=\binom mr=\binom{m+0}r$.

Die Behauptung gelte für $n$ und alle $r$. Für $r=0$ ist unmittelbar
$V_{n+1}(0)=1=\binom{m+n+1}0$. Sei nun $r\ge1$.
Wir spalten den letzten Summanden $k=r$ ab. In den übrigen Summanden
ist $r-k\ge1$, sodass Pascal ohne negative Indizes anwendbar ist:
\[
 \begin{aligned}
 V_{n+1}(r)
 &=\binom mr+\sum_{k=0}^{r-1}\binom mk\binom{n+1}{r-k}\\
 &=\binom mr+\sum_{k=0}^{r-1}\binom mk
       \left[\binom n{r-k}+\binom n{r-1-k}\right]\\
 &=V_n(r)+V_n(r-1).
 \end{aligned}
\]
Die erste Teilsumme wird durch den abgespaltenen Term
$\binom mr=\binom mr\binom n0$ zu $V_n(r)$ ergänzt;
die zweite ist bereits $V_n(r-1)$.
Mit der Induktionsannahme bei $r$ und $r-1$ und nochmals Pascal folgt
\[
 V_{n+1}(r)=\binom{m+n}r+\binom{m+n}{r-1}
          =\binom{m+n+1}r.
\]
Die Annahme musste für \emph{alle} $r$ gelten, weil beide Nachbarwerte
benötigt werden. Damit ist der Schritt $n\to n+1$ vollständig.
\end{proof}

\subsection*{Eine Summe von Quadraten}

\begin{satz}[Quadratsumme einer Zeile]
Für jedes $n\in\N$ gilt
\begin{equation}\tag{Q}\label{bin:squares}
 \sum_{k=0}^n\binom nk^2=\binom{2n}n.
\end{equation}
\end{satz}
\Beweisidee Schneide jeden Weg nach $(2n,n)$ in der Mitte, nach $n$
Schritten. Zur Station $(n,k)$ gibt es $\binom nk$ Hinwege; von dort
noch $\binom n{n-k}=\binom nk$ Weiterwege. Deshalb entsteht an jeder
Station ein Quadrat. Die Formel folgt genau dieser Zerlegung:
\begin{proof}
Setze in \eqref{bin:vandermonde} die beiden oberen Indizes gleich $n$
und den unteren Index ebenfalls gleich $n$. Dann folgt
\[
 \binom{2n}n=\sum_{k=0}^n\binom nk\binom n{n-k}.
\]
Für jeden Summanden ist $0\le k\le n$. Nach der Symmetrie
\eqref{bin:symmetry} ist deshalb $\binom n{n-k}=\binom nk$.
Das macht jedes Produkt zum behaupteten Quadrat.
\end{proof}

\begin{beispiel}
Die Zeile $3$ ist $1,3,3,1$. Ihre Quadratsumme beträgt
\[
 1^2+3^2+3^2+1^2=20=\binom63.
\]
Die gewöhnliche Zeilensumme wäre dagegen $1+3+3+1=8=2^3$.
Bei $n=0$ ergeben beide Formeln den Wert $1$.
\end{beispiel}

\section{Zusammenhang und Zuordnung zum Leitfaden}
\label{bin:map}

Alle Wegeideen beruhen auf derselben bewiesenen Zählinterpretation.
Wir spiegeln Wege, ordnen sie nach Endpunkten oder dem letzten
Rechtsschritt und zerlegen sie an einer Zwischenzeile. Die arithmetischen
Beweise sichern dieselben Zusammenhänge mit Rekursion und Induktion.

Die lokalen Gleichungsmarken (P) bis (Q) dienen der Orientierung.
Für die Zuordnung zum Leitfaden gelten die folgenden Nummern:

\begin{center}
\small
\renewcommand{\arraystretch}{1.25}
\begin{tabularx}{\linewidth}{@{}p{.35\linewidth}X@{}}
\toprule
Leitfaden Band 10 & Inhalt und Stelle\\
\midrule
\GuideRef{Def.75}{Definition 10H.5.3.2}
 & Pascal-Rekursion, Existenz und Eindeutigkeit
 (Abschnitt~\ref{bin:construction})\\
\GuideRef{Thm.76}{Theorem 10H.5.3.2}
 & Randwerte, Nullen, $\binom n1=n$, Symmetrie, Fakultätenformel
 (Abschnitte~\ref{bin:basic} und~\ref{bin:factorials})\\
\GuideRef{Thm.77}{Theorem 10H.5.3.3}
 & Binomischer Lehrsatz
 (Abschnitt~\ref{bin:theorem})\\
\GuideRef{Thm.78}{Theorem 10H.5.3.4}
 & Zeilen- und Diagonalensummen
 (Abschnitt~\ref{bin:sums})\\
\GuideRef{Thm.79}{Theorem 10H.5.3.5}
 & Vandermonde und Quadratsumme
 (Abschnitt~\ref{bin:vandermonde-section})\\
\bottomrule
\end{tabularx}
\end{center}

\textbf{Weiterlesen.}
Die \NewtonLink{} entwickelt Folgen aus ihren Anfangsdifferenzen.
Die Pascal-Rekursion und die Diagonalensummen bereiten diesen späteren
Einsatz vor. 
